本实验在本地编写攻击程序，并通过访问实验指导书给出的解密服务器完成 Padding Oracle 攻击。实验过程中，攻击者不需要知道 AES 密钥，只需依据服务器返回的状态判断解密结果的 PKCS\#7 填充是否合法。

本小组成员学号末位分别为 2、7、6，因此目标密文编号按下式确定：
\[
    2 + 7 + 6 = 15 \equiv 5 \pmod{10}
\]
故本实验选择 5 号目标密文进行攻击。

\begin{table}[H]
    \centering
    \caption{实验环境}
    \begin{tabular}{|c|c|}
        \hline
        项目 & 配置或说明 \\
        \hline
        操作系统 & Windows 11 \\
        \hline
        编程语言 & Python 3 \\
        \hline
        开发工具 & Visual Studio Code \\
        \hline
        加密算法 & AES-128-CBC \\
        \hline
        填充标准 & PKCS\#7 \\
        \hline
        目标密文编号 & 5 号密文 \\
        \hline
        密钥组 & Key\_2 \\
        \hline
        解密接口 & \verb|http://10.102.33.67:8208/dec_2?data=| \\
        \hline
    \end{tabular}
\end{table}

实验指导书中的解密服务器会对输入密文进行 AES-128-CBC 解密，并检查解密后明文是否满足 PKCS\#7 填充规则。若填充正确，服务器返回 HTTP 200；若填充错误，则返回 HTTP 500。因此，本实验通过构造查询密文并观察服务器反馈状态，逐步恢复目标密文对应的明文。

实验材料中给出的 5 号密文属于 Key\_2 加密所得密文，因此程序访问的解密接口为 \verb|dec_2|。实验运行时需要保证本机能够访问学校实验服务器，通常需要处于校园网环境或连接学校 VPN。
